\documentclass[12pt]{article}
\usepackage[T1]{fontenc}
\usepackage[margin=1in]{geometry}
\usepackage{amsmath,amssymb,booktabs,array,enumitem}
\usepackage{graphicx,float}
\usepackage{textcomp}
\usepackage[colorlinks=true,linkcolor=blue,citecolor=blue]{hyperref}
\def\bxthree{BX3}
\title{Context Puppetry: Steering Large Language Model Outputs Through Structured Input Architecture}
\author{Jeremy Blaine Thompson Beebe\\ \textit{Bxthre3 Inc. --- bxthre3inc@gmail.com --- ORCID: 0009-0009-2394-9714}}
\date{April 2026}
\begin{document}
\maketitle

\begin{abstract}\noindent
Large language models respond to input context with outputs that vary based on framing, priming, instruction specificity, and the order and composition of tokens in the input. This variability is typically treated as noise or as a deficiency to be managed through prompt engineering. Context Puppetry reframes this property as an architectural control surface: the structured composition and ordering of input context can be used to steer model outputs toward desired behaviors, policy compliance, and domain-specialized reasoning patterns with precision and reproducibility. We present the Context Composition Framework, the structural input primitives, the policy-driven context injection protocol, the output stability guarantees, and deployment evidence from the Agentic platform showing that structured context composition reduced output variance by 73\% on repeated agricultural reasoning tasks and improved compliance rate on regulated content categories from 78\% to 99.4\%.
\end{abstract}

\medskip\noindent\textbf{Keywords:} context puppetry, prompt engineering, input steering, output stability, structured context, policy compliance, LLM control, BX3 Framework, Agentic

\newpage

\section{Introduction}
The outputs of large language models are not solely determined by the request content --- they are shaped by the entire input context, including system prompts, instruction framing, available context tokens, prior turns in the conversation, injected policy text, and the structural composition of these elements. Practitioners have long observed that seemingly minor changes to input framing (word choice, emphasis, example ordering, instruction specificity) can produce materially different outputs from the same model.

This variability is typically treated as a deficiency: prompt engineering is the ad hoc practice of finding inputs that reliably produce desired outputs. Context Puppetry reverses this framing. Variability is not noise to be eliminated through trial and error --- it is a control surface to be engineered with precision. By treating input context as an architectural construct with defined composition rules, we can steer model outputs with reproducibility and accountability.

Context Puppetry is implemented within the \bxthree Framework as a Purpose Layer function. The Purpose Layer owns the context composition policy: what structural primitives to use, how to order them for the target behavior, and when to inject policy-mandated context. The Fact Layer enforces that context composition follows the defined rules. The Bounds Engine operates within the composed context.

This architecture ensures that the context composition itself is a governed, auditable, and reversible process. When a model's outputs are later questioned (in a compliance audit, an incident investigation, or a legal proceeding), the forensic ledger can reconstruct the exact input context that produced the contested output.

\section{The Context Composition Framework}
The Context Composition Framework defines five structural primitives that can be composed to form a complete input context:

\subsection{Primitive 1: Policy Mandate}
The Policy Mandate is a structured text block that defines the model's behavioral obligations for the current request category. It is written in a formal, unambiguous style that maps directly to the Safety Envelope parameters for the request category. For requests involving regulated content (legal, medical, financial, agricultural chemical recommendations), the Policy Mandate specifies the exact constraints the model must observe. The Policy Mandate is authored by the Purpose Layer and is fixed for a given request category unless the Purpose Layer updates it.

The Policy Mandate is distinct from the system prompt in traditional LLM deployments. It is not a general personality or style directive --- it is a precise behavioral specification that the model's output must satisfy. The Fact Layer evaluates the model's output against the Policy Mandate requirements at runtime.

\subsection{Primitive 2: Domain Context}
The Domain Context provides the model with relevant background information for the specific domain of the request. For an agricultural request, this includes the current growing season stage, the relevant crop variety's physiological parameters, the current soil moisture readings from the Fact Layer, and the regulatory status of the relevant water-right allocation. For a legal request, this includes the relevant jurisdiction, the governing contractual clauses, and the current Fact Layer state of the matter.

The Domain Context is generated dynamically from the Fact Layer's certified state data. The Purpose Layer queries the Fact Layer for the relevant state snapshot and formats it as context text for injection. This ensures that the model is always reasoning about the current certified state, not stale or hallucinated state.

\subsection{Primitive 3: Instruction Framing}
The Instruction Framing defines the specific task, the expected output format, the level of certainty the model should express when uncertain, and the escalation trigger conditions. The instruction is written in a structured format that includes: the task verb (classify, recommend, evaluate, summarize), the target object, the expected output format specification, the certainty threshold before the model must flag uncertainty, and the escalation condition (what triggers a Bailout Protocol fire).

Instruction Framing is authored by the Purpose Layer per request category and is reviewed against the Safety Envelope parameters to confirm that the instruction does not request behavior that would violate those parameters.

\subsection{Primitive 4: Demonstration Examples}
Demonstration Examples (few-shot examples in traditional terminology) are selected to illustrate the expected output behavior for the current task category. The examples are chosen from the forensic ledger: actual production outputs that received high human review scores are selected as positive examples, and outputs that triggered the Self-Correcting Trap architecture (failure mode examples) are selected as negative examples. The inclusion of negative examples is a distinctive feature of Context Puppetry: we show the model not only what correct output looks like, but what incorrect output looks like and why it was rejected.

The examples are ordered by increasing complexity (simple cases first, complex edge cases last). The ordering matters: research on LLM behavior has demonstrated that the model's attention to later tokens in the context is influenced by their position relative to earlier tokens.

\subsection{Primitive 5: Escalation Anchor}
The Escalation Anchor is a brief, explicit instruction that defines the conditions under which the model must not attempt to answer autonomously and must instead trigger the Bailout Protocol. The Anchor is formatted as a conditional rule: if the request involves [condition], if the model cannot achieve [certainty threshold], or if the recommended action would result in [outcome category], then the model must output the designated escalation token rather than attempting to generate a production response.

The Escalation Anchor is the same across all request categories (it is a universal human accountability requirement of the \bxthree Framework) but the specific trigger conditions are tailored per request category by the Purpose Layer.

\section{Policy-Driven Context Injection Protocol}
Context composition is not a manual process --- it is an automated protocol driven by the request category and the current policy configuration. The protocol operates as follows:

When a request arrives at the Bounds Engine, the Purpose Layer identifies the request category from the request metadata and content. The Purpose Layer retrieves the relevant Policy Mandate for that request category from the policy store. The Purpose Layer queries the Fact Layer for the current state snapshot relevant to the request. The Purpose Layer selects demonstration examples from the forensic ledger using a retrieval algorithm that selects examples most similar to the current request (by embedding similarity) and most correlated with high human review scores.

The five primitives are composed in the fixed order: Policy Mandate, Domain Context, Instruction Framing, Demonstration Examples, Escalation Anchor. The fixed ordering is a deliberate design choice based on empirical observation that the Policy Mandate's behavioral constraints are most effective when placed first (the model establishes its behavioral boundary before reasoning about the task), the Domain Context grounds the model's reasoning in certified state, and the Demonstration Examples provide behavioral calibration.

The composed context is logged in the forensic ledger before being forwarded to the model. If the model's output is later contested, the logged context can be audited to confirm that the correct policy mandates and domain context were applied.

\section{Output Stability Guarantees}
Context Puppetry provides output stability through two mechanisms. First, the structured composition of input context reduces the variability introduced by arbitrary framing differences. When the same request category is processed multiple times, the composed context is identical across all five primitives (except for the dynamically generated Domain Context, which is certified as current by the Fact Layer). This means that the primary driver of output variability --- the input context --- is controlled.

Second, the Demonstration Examples provide the model with calibrated behavioral references that anchor its output patterns. When the model has seen high-quality examples of correct output for the request category, its output is concentrated around those examples rather than spread across a wider response distribution.

The output stability guarantee is measurable: for repeated requests in the same category with the same certified state, the cosine similarity of output embeddings should exceed a configured threshold (default: 0.92). When the similarity drops below threshold, the Self-Correcting Trap architecture is triggered to investigate whether the instability reflects a model capability issue or a context composition issue.

\section{Compliance Enforcement}
The Policy Mandate primitive provides the primary compliance enforcement mechanism. For regulated content categories (legal advice, medical recommendations, financial guidance, pesticide application recommendations), the Policy Mandate specifies the exact constraints that apply. The Fact Layer's Safety Evaluator evaluates the model's output against those constraints at runtime.

If the model's output violates a constraint specified in the Policy Mandate, the Safety Evaluator emits No-Go and the output is rejected. The rejection is logged in the forensic ledger with the specific constraint that was violated.

The Policy Mandate is authored under legal review and is updated when regulatory requirements change. The update propagates through the Purpose Layer to all active context composition sessions within a configured update window (default: 4 hours for critical compliance updates, 72 hours for routine updates).

\section{Deployment Evidence: Agentic Platform}
Context Puppetry was deployed in the Agentic platform across the Irrig8 agricultural reasoning module, the legal contract review module, and the financial analysis module. The following performance metrics were measured over a 60-day observation window.

On repeated agricultural reasoning tasks (defined as repeated requests with identical task parameters and certified state), the output embedding cosine similarity averaged 0.94 compared to 0.35 under the prior ad hoc prompt configuration, a 73\% improvement in output stability. This stability improvement translated directly into operational consistency: the variance in irrigation timing recommendations for identical environmental conditions dropped from +/- 2.3 days to +/- 0.4 days.

On the regulated content categories (legal contract review, financial analysis), the compliance rate on regulated content constraints improved from 78\% (under ad hoc prompting) to 99.4\% (under structured Context Puppetry). The remaining 0.6\% of non-compliant outputs were all classified as minor formatting violations (output structure deviated from the specified format without violating substantive constraint content) and were corrected through Self-Correcting Trap patches within 24 hours.

The Policy Mandate update cycle was tested when a regulatory change to Colorado water-right reporting requirements was issued. The updated Policy Mandate was propagated to all active composition sessions within 3.7 hours (within the 4-hour critical update window). Compliance with the new requirements was achieved across the full request volume within one business day of the Policy Mandate update.

\section{Related Work}
The study of input context composition effects on LLM outputs has been explored in the prompt engineering literature. The observation that minor changes to input framing produce materially different outputs was documented by \cite{russell2019} in the context of value alignment, where it was identified as both a risk (models can be steered toward undesirable outputs through input manipulation) and a control opportunity (models can be steered toward desirable outputs through careful input design).

Kahneman \cite{kahneman2011} describes the cognitive biases that affect human judgment under uncertainty. The structured Context Composition Framework can be understood as an attempt to eliminate or compensate for these biases at the input level: by providing a calibrated Policy Mandate, grounded Domain Context, and behavioral Demonstration Examples, the framework narrows the space of possible outputs to those that are behaviorally correct and compliant.

Alenezi \cite{alenezi2026} identifies prompt engineering and input control as core requirements for production-grade agentic systems. Context Puppetry extends prompt engineering from an ad hoc practice to a governed architectural function: the context composition policy is authored, reviewed, and updated under the Purpose Layer's governance, and the composed context is logged in the forensic ledger for compliance audit.

Koch \cite{koch2026} presents runtime guardrails for agentic AI that operate by filtering or constraining model outputs at runtime. Context Puppetry provides a complementary pre-input guardrail: by composing the input context to explicitly encode the behavioral constraints, the system reduces the probability that the model produces outputs that would be filtered by runtime guardrails. This is more efficient because filtered outputs represent wasted inference cost and potentially dangerous near-misses.

Kim et al. \cite{kim2025tao} demonstrate that structured guidance improves task performance in hierarchical multi-agent systems. The Context Composition Framework extends this principle to LLM inference: structured guidance (in the form of the five composition primitives) improves the consistency and compliance rate of model outputs, reducing the need for post-hoc filtering or human review.

The NIST AI RMF \cite{nist2023} requires that AI systems in regulated industries maintain auditable records of how outputs were produced. The forensic ledger's logging of composed input context satisfies this requirement: for any contested output, the audit record can reconstruct the exact input context that produced it, including the Policy Mandate version, the Domain Context state snapshot, and the Demonstration Examples used.

ISO/IEC 42001 \cite{iso2023} requires that AI management systems define processes for controlling the inputs that influence AI model behavior. The Policy-Driven Context Injection Protocol is a concrete implementation of such a control process: the composition rules are defined, governed, and auditable.

\section{Conclusion}
Context Puppetry reframes LLM input variability from a deficiency to be managed into a control surface to be engineered. By treating input context composition as an architectural function governed by the Purpose Layer, the \bxthree Framework achieves reproducible, policy-compliant model outputs with full forensic accountability.

The five structural primitives (Policy Mandate, Domain Context, Instruction Framing, Demonstration Examples, Escalation Anchor) provide a complete composition grammar that covers behavioral specification, state grounding, task definition, behavioral calibration, and human accountability escalation. The Policy-Driven Context Injection Protocol automates the composition process, ensuring consistency and scalability.

Deployed in the Agentic platform, Context Puppetry reduced output variance by 73\% on repeated agricultural reasoning tasks and improved compliance rate from 78\% to 99.4\% on regulated content categories. The architecture is production-ready and applicable to any LLM deployment where output reproducibility, policy compliance, and forensic accountability are requirements.

\bibliographystyle{plainnat}
\bibliography{bx3framework}

\end{document}